LLM-guided theorem provers achieve 89\% success on miniF2F olympiad problems while pure automated provers manage only 16\% --- a 50-percentage-point gap dominating recent literature. Yet no prior work quantifies \emph{which mechanisms} drive this advantage. We hypothesize the gap decomposes into three testable mechanisms: natural language understanding (60\% contribution), proof depth capability (30\%), and corpus pattern matching (10\%). Through controlled ablation studies on miniF2F Lean 4 (N=244), we establish: (1) pure automated prover (lean-auto) baseline at 15.6\% [11.5\%, 20.3\%], (2) removing natural language hints drops LLM success from 62.3\% to 32.8\% ($\Delta$=29.51\% [20.90\%, 38.11\%], p<10$^{-9}$), validating $\sim$60\% contribution, (3) proof depth filtering shows only 3.7\% effect (below 5\% threshold), rejecting the 30\% hypothesis. Natural language understanding is the \textbf{dominant mechanism} --- LLMs parse informal hints (``for all prime p'') into formal tactics (\texttt{Nat.Prime} lemmas) that automated provers ignore. This finding shifts design priorities from architectural scaling (longer context for deep proofs --- rejected) to NL-formal integration (hint extraction, semantic tactic libraries). For hybrid systems, mechanistic attribution suggests routing NL-rich problems to LLMs and formal-only goals to automated provers, explaining Thor's 8.2\% hybrid synergy. Our controlled ablation framework --- with falsification thresholds enforcing rigor --- provides the first quantified mechanistic attribution of LLM theorem proving advantage.
